\chapter{Expressions, identités et formules}
\section*{Exercice 5.1 -- Développer}
Développer $(2x-3)^2$ et $(x+4)(x-4)$.

\section*{Exercice 5.2 -- Factoriser}
Factoriser $6x^2-15x$ et $(x+1)(2x-3)+4(x+1)$.

\section*{Exercice 5.3 -- Domaine conservé}
Simplifier $(x^2-9)/(x+3)$ en précisant son domaine. Peut-on attribuer une valeur à l'expression initiale en $-3$ ?

\section*{Exercice 5.4 -- Isoler}
Dans $V=\pi r^2h$, exprimer $h$ si $r>0$, puis $r$ si $h>0$, $V\geq0$ et $r\geq0$.

\section*{Exercice 5.5 -- Choisir la forme}
Écrire $x^2-8x+20$ sous la forme d'un carré plus une constante. En déduire son minimum sur $\R$ et résoudre l'égalité à $8$.

\section*{Exercice 5.6 -- Moyennes}
Approfondissement : pour $a,b>0$ de somme $12$, montrer $ab\leq36$ et trouver le cas d'égalité.
